\addtocontents{toc}{\protect\newpage}

\section{Manual de usuario} \label{sec:appendix_A_user_manual}

El presente apéndice describe el flujo básico de uso de la aplicación
desarrollada. Su propósito es servir como
guía operativa para la puesta en marcha del entorno, la ejecución de la
aplicación, la interacción con la interfaz durante una sesión completa de
tratamiento de imagen y la exportación del perfil de operaciones que actúa como
entrada del módulo de \textit{benchmarking}.

\subsection{Preparación del entorno}

La aplicación interactiva se ejecuta como una aplicación de escritorio basada en
Python y PyQt6. Como paso previo, debe prepararse un entorno virtual e instalarse
las dependencias del proyecto.

\subsubsection{Creación del entorno virtual}

El \autoref{lst:venv_setup} muestra los comandos necesarios para crear y activar un
entorno virtual desde el directorio \texttt{app/} del proyecto.

\begin{lstlisting}[language=bash, caption={Creación y activación del entorno virtual}, label={lst:venv_setup}]
cd app
python -m venv .venv
source .venv/bin/activate
\end{lstlisting}

\subsubsection{Instalación de dependencias}

Una vez activado el entorno, deben instalarse las dependencias definidas
en \texttt{requirements.txt}, como se muestra en el \autoref{lst:install_dependencies}.

\begin{lstlisting}[language=bash, caption={Instalación de dependencias de la aplicación}, label={lst:install_dependencies}]
pip install --upgrade pip
pip install -r requirements.txt
\end{lstlisting}

\noindent Debe tenerse en cuenta que el uso de GPU depende de la disponibilidad
de CUDA en el sistema y de una instalación compatible de PyTorch. En ausencia de
soporte CUDA, la aplicación permanece operativa utilizando CPU. La instalación 
de CUDA puede realizarse siguiendo las instrucciones oficiales de NVIDIA, 
disponibles en \url{https://developer.nvidia.com/cuda-downloads}.

\newpage

\subsection{Puesta en marcha de la aplicación}

La aplicación puede ejecutarse en modo normal o en modo de desarrollo con
\textit{hot reload}. Para la utilización habitual del sistema, basta con lanzar
la aplicación principal.

\subsubsection{Ejecución estándar}
En el \autoref{lst:launch_app} se muestra el comando para iniciar la aplicación interactiva
desde el directorio \texttt{app/}.

\begin{lstlisting}[language=bash, caption={Lanzamiento de la aplicación interactiva}, label={lst:launch_app}]
cd app
source .venv/bin/activate
python app.py
\end{lstlisting}

\subsection{Visión general de la interfaz}

Al iniciarse la aplicación, se muestra la ventana principal maximizada, como se observa 
en la \autoref{fig:appendix_a_gui_overview}.

\begin{figure}[h!]
    \centering
    \includegraphics[width=0.9\textwidth]{sections/appendix/A/img/1.png}
    \caption{Visión general de la aplicación interactiva}
    \label{fig:appendix_a_gui_overview}
\end{figure}

\newpage

\subsection{Flujo básico de uso}

Se describe a continuación una sesión típica de trabajo, desde la carga de una
imagen hasta la generación del perfil de operaciones reutilizable por el módulo
de \textit{benchmarking}.

\subsubsection{Carga de imagen}

El primer paso consiste en pulsar el botón \textit{Open file...} del panel
izquierdo y seleccionar una imagen compatible con alguno de los formatos
soportados por la aplicación. Tras la carga, el sistema actualiza la vista
principal, la información de la imagen y los elementos de análisis visual
disponibles, como puede observarse en la \autoref{fig:appendix_a_load_image}.

\begin{figure}[h!]
    \centering        
    \includegraphics[width=0.9\textwidth]{sections/appendix/A/img/2.png}
    \caption{Carga inicial de una imagen en la aplicación}
    \label{fig:appendix_a_load_image}
\end{figure}

Una vez cargada la imagen, la aplicación muestra en el panel izquierdo un
histograma RGB y un gráfico de perfiles centrales de intensidad, mientras que el
visor central refleja tanto la imagen como la información técnica asociada.

\subsubsection{Inspección visual inicial}

Antes de modificar el \textit{pipeline}, resulta recomendable realizar una
inspección visual inicial de la imagen cargada. Para ello, el usuario puede:

\begin{itemize}[noitemsep]
    \item observar la distribución tonal mediante el histograma RGB,
    \item analizar los perfiles centrales de intensidad,
    \item revisar los metadatos mostrados en la pestaña \textit{Image Info},
    \item utilizar zoom y desplazamiento sobre el visor principal.
\end{itemize}

\newpage

\subsubsection{Definición del \textit{pipeline} de procesado}

El núcleo de interacción de la herramienta es la construcción del
\textit{pipeline} de operaciones. Desde el selector del panel izquierdo, el
usuario puede añadir nuevas operaciones a la secuencia de procesado. Cada nueva
operación se incorpora al final del \textit{pipeline} y aparece, a su vez, en el
panel derecho con su control de parámetros correspondiente. En la 
\autoref{fig:appendix_a_pipeline_build} se muestra el resultado tras añadir una 
operación.

\begin{figure}[h!]
    \centering
    \includegraphics[width=0.9\textwidth]{sections/appendix/A/img/3.png}
    \caption{Construcción inicial del \textit{pipeline} de procesado}
    \label{fig:appendix_a_pipeline_build}
\end{figure}

El flujo normal de uso consiste en encadenar operaciones de forma
incremental, observando tras cada modificación el efecto resultante sobre la
imagen procesada y ajustando los parámetros de cada operación según sea necesario.

\newpage

\subsubsection{Ajuste de parámetros}

Cada operación activa presenta en el panel derecho un conjunto de controles
específico. Dichos controles permiten modificar parámetros numéricos, seleccionar
métodos o ajustar configuraciones particulares según la naturaleza de cada
operación, como se muestra en la \autoref{fig:appendix_a_parameters}.

\begin{figure}[h!]
    \centering
    \includegraphics[width=0.76\textwidth]{sections/appendix/A/img/4.png}
    \caption{Modificación de parámetros de las operaciones activas}
    \label{fig:appendix_a_parameters}
\end{figure}

Durante esta fase, el usuario puede reajustar tantas veces como resulte necesario
la secuencia de operaciones, observando en tiempo real su efecto sobre la imagen.

\subsubsection{Guardado de la imagen procesada}

Una vez alcanzado un resultado visual satisfactorio, la imagen procesada puede
guardarse pulsando el botón \textit{Save file...}; esta operación genera una
salida en formato \texttt{.jpg} a partir del estado actual del \textit{pipeline}.
En la \autoref{fig:appendix_a_save_image} se muestra el diálogo de guardado.

\begin{figure}[h!]
    \centering
    \includegraphics[width=0.76\textwidth]{sections/appendix/A/img/5.png}
    \caption{Guardado de la imagen procesada}
    \label{fig:appendix_a_save_image}
\end{figure}

\newpage

\subsection{Generación del perfil de operaciones}

Uno de los resultados más relevantes de la interacción con la aplicación es la
posibilidad de exportar el \textit{pipeline} de operaciones en formato
\textit{json}.

Para ello, debe utilizarse el botón \textit{Generate profile} del panel izquierdo.
La aplicación abrirá un diálogo de guardado y generará un fichero
\texttt{operations\_profile.json} o cualquier otro nombre especificado por el
usuario. En la \autoref{fig:appendix_a_export_profile} se muestra 
el diálogo de exportación del perfil.

\begin{figure}[h!]
    \centering
    \includegraphics[width=0.8\textwidth]{sections/appendix/A/img/6.png}
    \caption{Exportación del perfil de operaciones del \textit{pipeline}}
    \label{fig:appendix_a_export_profile}
\end{figure}

Conceptualmente, dicho fichero contiene una lista ordenada de operaciones, donde
cada entrada almacena el nombre de la operación y el conjunto de parámetros
asociados a ella. Este fichero sigue la estructura mostrada en el \autoref{lst:profile_structure}.

\begin{lstlisting}[caption={Estructura conceptual de un perfil exportado}, label={lst:profile_structure}]
[
    {
        "operation": "Debayer",
        "params": {
            "algorithm_name": "debayer5x5"
        }
    },
    {
        "operation": "MinMaxPercentileNormalization",
        "params": {
            "lower_percentile": 0.01,
            "upper_percentile": 0.98
        }
    },
    ...
]
\end{lstlisting}

\newpage

\subsection{Ejecución del \textit{benchmarking} a partir del perfil exportado}

Una vez generado el perfil, este puede emplearse como entrada directa del módulo
de \textit{benchmarking}. La ejecución se realiza desde el directorio
\texttt{app/} mediante \textit{pytest}, utilizando el argumento
\texttt{-{}-bench-profile} para indicar la ruta del fichero \textit{json}, como
se muestra en el \autoref{lst:benchmark_execution}.

\begin{lstlisting}[language=bash, caption={Ejecución del benchmarking sobre un perfil exportado}, label={lst:benchmark_execution}]
cd app
source .venv/bin/activate
pytest benchmark/run.py --bench-profile=operations_profile.json
\end{lstlisting}

Si el perfil se encuentra almacenado en otra ubicación o bajo otro nombre, basta
con sustituir la ruta anterior por la correspondiente. Se muestra un ejemplo en el \autoref{lst:benchmark_execution_explicit}.

\begin{lstlisting}[language=bash, caption={Ejemplo de ejecución con ruta explícita}, label={lst:benchmark_execution_explicit}]
pytest benchmark/run.py --bench-profile=pipelines/pipeline_1.json
\end{lstlisting}

Durante la ejecución, el módulo construye dinámicamente el \textit{pipeline}
definido en el perfil y lo evalúa sobre distintas combinaciones de tamaño de
imagen, tipo de dato y dispositivo de procesamiento. Los resultados se almacenan 
automáticamente en el directorio \texttt{.benchmarks/}. Desde la terminal se muestra un 
resumen de los resultados obtenidos, como se puede observar en la \autoref{fig:appendix_a_benchmark_run}.

\begin{figure}[h!]
    \centering
    \includegraphics[width=\textwidth]{sections/appendix/A/img/7.png}
    \caption{Ejemplo de salida en terminal tras la ejecución del \textit{benchmarking}}
    \label{fig:appendix_a_benchmark_run}
\end{figure}
